\documentclass[10pt,a4paper]{amsart}
\usepackage[margin=19mm]{geometry}
\usepackage{amsmath,amssymb,mathtools}
\usepackage[scr=rsfs,cal=euler]{mathalfa}
\usepackage{xfrac}
\usepackage[unicode,hidelinks,bookmarks=false]{hyperref}
\newcommand{\ie}{i.e.\ }
\newcommand{\cf}{cf.\ }
\newcommand{\resp}{respectively\ }
\newcommand{\Subsection}{\subsection}
\providecommand{\nobreakdash}{\nobreak\hbox{-}}
\setlength{\emergencystretch}{2em}
\multlinegap0pt
\usepackage{amssymb}
\usepackage{enumerate}
\usepackage{xcolor}
\usepackage{mathtools}
\usepackage[shortlabels]{enumitem}
\newtheorem{thm}{Theorem}
\newtheorem{cor}{Corollary}
\newcommand{\RR}{\mathbb{R}}
\title{On the number of pairwise touching cylinders in $\mathbb{R}^d$}
\author{Jozsef Solymosi \and Joshua Zahl}
\date{Source: arXiv:2512.24595v1 (CC BY 4.0)}
\begin{document}
\maketitle

\noindent\emph{Source and licence.} On the number of pairwise touching cylinders in $\mathbb{R}^d$. Version arXiv:2512.24595v1 (CC BY 4.0), distributed by arXiv under the Creative Commons Attribution 4.0 International licence, http://creativecommons.org/licenses/by/4.0/, as recorded in the arXiv licence metadata for this exact version and shown on the abstract page https://arxiv.org/abs/2512.24595v1. Full licence notice: \texttt{licenses/arxiv-2512.24595-CC-BY-4.0.txt}.\par\smallskip

\noindent\textit{Editorial scope.} Counted unit: the article's thm unitDistanceCylindersLines (source0.tex lines 74--76). The article's complete original proof of it is delivered with it (source0.tex lines 187--193, one \texttt{proof} environment of 7 lines, which the article places away from the statement). No mathematics is written, repaired or re-derived: the statement and the proof are the article's own lines, in the article's own notation, and no retained line is substituted or re-flowed.  Retained uncounted as the source-local context the proof needs: lines 71--73; lines 158--160; lines 167--169.  Nothing was omitted from any retained span, so the package carries no omission record.  No retained line contains a citation, so the wrapper carries no bibliography and no bibliography entry.  No retained line contains a figure environment, an image-inclusion command or drawing code, so no image is delivered and the package declares no asset dependency.  Editorial formatting only, in addition: an AMS \texttt{amsart} preamble that restates the article's own declarations and macros used by the retained lines, namely the environments \texttt{thm}, \texttt{cor} on separate counters, together with the standard packages those lines need.  The retained environments print in this document as Theorem 1, Corollary 1 in order of appearance, whereas the article prints the same items with the numbering of its own full document.  The complete native source of the article as distributed in the arXiv e-print for this exact version is bundled unchanged as \texttt{sources/191913/source0.tex}.
\begin{thm}\label{unitDistanceLines}
Let $C$ be a set of mutually touching unit cylinders in $\RR^d$. Then $|C|\leq 4d\cdot 7^{2d-3}$.
\end{thm}

\begin{cor}\label{consequenceOfMilnorThom}
Let $X\subset\RR^{2d-2}$. Then $\overline X^F$ is a union of at most $4\cdot 7^{2d-3}$ Euclidean connected components.%, and $\overline Y^G$ is a union of at most $4\cdot 7^{2d-2}$ Euclidean connected components.
\end{cor}

\begin{cor}\label{consequenceOfBPR}
Let $Y\subset\RR^{2d-1}$. Then $\overline Y^G \cap \{(x,r)\in\RR^{2d-1}\times\RR\colon r \geq 1\}$ is a union of at most $20\cdot 7^{2d-2}$ Euclidean connected components.
\end{cor}

\begin{thm}\label{unitDistanceCylindersLines}
Let $C$ be a set of mutually touching cylinders (with arbitrary radii) in $\RR^d$. Then $|C|\leq 20(d+1)\cdot 7^{2d-2}$.
\end{thm}

\begin{proof}[Proof of Theorem \ref{unitDistanceCylindersLines}]
The proof proceeds in a similar fashion to that of Theorem \ref{unitDistanceLines}. Since at most $d+1$ circles in $\RR^d$ can mutually touch, it suffices to prove that every set of mutually touching cylinders with distinct directions has cardinality at most $20\cdot 7^{2d-2}$. Suppose for contradiction that there exists such a set of cylinders with cardinality $20\cdot 7^{2d-2}+1$. After rescaling we may suppose that each cylinder has radius $\geq 1$. Let $Y\subset\RR^{2d-1}$ be the corresponding set of points (the final coordinate denotes the radius of the cylinder). We wish to obtain a contradiction. 

Arguing as above and using Corollary \ref{consequenceOfBPR} in place of Corollary \ref{consequenceOfMilnorThom}, we obtain a Euclidean connected set $U\subset \bar Y^G \cap \{(x,r)\in\RR^{2d-1}\times\RR\colon r \geq 1\}$ that contains at least two points of $Y$, and thus $\pi(U)$ is uncountable. % (as above, $\pi\colon\RR^{2d-1}\to\RR^{d-1}$ denotes the projection to the first $d-1$ coordinates, corresponding to the direction of the tube). 
%
On the other hand, if $(x,r),\ (y,s)\in U$ and $\pi(x,r)\neq\pi(y,s)$, then $\Vert x-y\Vert \geq \operatorname{dist}(L_x, L_y) \geq (r+s)\geq 2$. This means that $\pi(U)$ must be (at most) countable. This is our desired contradiction.
\end{proof}

\end{document}
